\chapter*{前言}
\addcontentsline{toc}{chapter}{前言}
\markboth{前言}{前言}
\setchapterheader{前言}

在笔记本上运行一个 Python 程序，就能让机器人站起来、走起来。这件事听起来很直接，真正动手时，我们却会接连碰到许多问题：命令应该在哪个终端输入？包装到了哪个 Python 里？为什么模型加载成功了，按手柄还是没有反应？

这些问题往往来自一些平时容易忽略的基础概念。因此，本手册先介绍终端、Linux、Python、虚拟环境和网络通信，再把它们放到 Bumi 的部署过程中。我们希望读者既能按步骤完成操作，也能理解每条命令在做什么，遇到问题时知道从哪里查起。

手册分为两部分。第一部分是前置知识，介绍开发中会用到的基本工具与原理；第二部分是部署实践，依次完成环境配置、网络连接、手柄输入、程序启动和运动控制。编译与运行统一使用 Conda 环境 \code{bumi_deploy}。

初次接触 Linux 和 Python 的读者，建议从第一部分读起。已经熟悉这些工具的读者，可以先看第\,\ref{chap:conda}\,章中的环境配置，再进入部署实践。每个部署步骤都列出了操作方法、预期结果和常见问题，附录则提供按键和故障速查。

本手册使用 Ubuntu 22.04 x86\_64 笔记本，通过网线连接 Bumi 后背网口，手柄接收器插在笔记本 USB 口。程序在笔记本上读取手柄、执行 ONNX 推理，并通过 DDS 与运控板交换数据。

调试时请全程使用吊架，按顺序完成模式切换。停止程序前，先通过手柄停控并确认机器人得到稳定支撑。具体操作见部署实践和附录 B。
